% Style commun des lettres de motivation — Ismaila AMADOU DEH
\documentclass[11pt,a4paper]{article}
\usepackage[utf8]{inputenc}
\usepackage[T1]{fontenc}
\usepackage[english,french]{babel}
\usepackage{lmodern}
\usepackage[margin=2.2cm,top=1.8cm]{geometry}
\usepackage{xcolor}
\usepackage{parskip}
\usepackage[hidelinks]{hyperref}
\definecolor{accent}{HTML}{1F4E79}
\definecolor{gris}{HTML}{6B675F}
\pagestyle{empty}
% Coordonnées privées (e-mail, téléphone) : fichier local non versionné
\newcommand{\contact}{}
\IfFileExists{coordonnees.tex}{\input{coordonnees.tex}}{}
\newcommand{\entete}{%
  {\LARGE\bfseries\color{accent} Ismaila AMADOU DEH}\\[2pt]
  {\small\color{gris} Étudiant M2 Physique Fondamentale -- parcours Théories Quantiques, Université de Tours\\
  \contact\href{https://www.linkedin.com/in/ismaila-amadou-deh}{linkedin.com/in/ismaila-amadou-deh} \;·\;
  \href{https://github.com/IsmailaDEH/Ismaila}{github.com/IsmailaDEH/Ismaila}}\\[-4pt]
  {\color{accent}\rule{\linewidth}{0.6pt}}}
\newcommand{\destinataire}[2]{\null\hfill\begin{minipage}[t]{0.45\linewidth}#1\\[6pt]Tours, le \today\end{minipage}\par\medskip
  \textbf{Objet :} #2\par}
\newcommand{\signature}{\par\bigskip\hfill Ismaila AMADOU DEH}

\begin{document}
\entete
\destinataire{Framatome\\Service Recrutement\\Courbevoie}{Candidature au stage « Ingénieur Mécanique -- Analyses dynamiques de la cuve et de ses internes »}

Madame, Monsieur,

Étudiant en Master~2 Physique Fondamentale à l'Université de Tours, je vous propose ma candidature pour ce stage de fin d'études. Physicien de formation, je suis attiré par la simulation de systèmes réels et par les enjeux de sûreté de l'industrie nucléaire.

Mon parcours m'a donné de bonnes bases en \textbf{dynamique des systèmes mécaniques}. J'ai modélisé le pendule double (projet de Licence~3), système non linéaire où j'ai mesuré l'exposant de Lyapunov ($\lambda = 0{,}75~\mathrm{s^{-1}}$) et contrôlé la conservation de l'énergie à $2\times10^{-7}$ près. J'ai aussi comparé plusieurs schémas d'intégration (Euler, RK4, Verlet symplectique) sur des problèmes d'orbites, en étudiant leur stabilité et leur ordre de convergence, ainsi que la résolution d'EDP par différences finies.

Ce qui me caractérise est la démarche : écrire les équations, coder et \textbf{valider contre une solution de référence} avant d'interpréter, puis quantifier les incertitudes. Je pense que cette rigueur est précieuse pour des analyses dynamiques (modes propres, réponse vibratoire, sollicitations sismiques) qui doivent être justifiées.

Je maîtrise Python (NumPy, SciPy), Git et LaTeX, et je suis prêt à me former rapidement aux codes de calcul de structures que vous utilisez.

Je serais heureux de vous en dire plus lors d'un entretien. Veuillez agréer, Madame, Monsieur, l'expression de mes salutations distinguées.
\signature
\end{document}
